\section{Introduction}
\label{sec:intro}

Autonomous Scientific Discovery (ASD) refers to the use of Large Language Models (LLMs)~\citep{yang2024qwen2-5,internvl2,qwen2vl,Deepseek-r1} and robotics to independently perform scientific research without direct human intervention~\citep{yuan2025dolphin,yan2025surveyforge,gottweis2025towards,yamada2025ai,lu2024ai}. This approach holds transformative potential for accelerating the pace of discovery across various scientific domains. By automating tasks such as data analysis, hypothesis generation, experiment design, and result interpretation, automated systems~\citep{yuan2025dolphin, lu2024ai} can efficiently process vast amounts of information and uncover patterns or insights that may be difficult for human researchers to detect.

ASR, while promising, faces significant challenges in generating effective and novel proposals, as well as achieving closed-loop feedback for the experimental validation of these proposals:

\begin{itemize}
    \item First, generating proposals that are both effective and novel is a complicated task. Autonomous systems must identify research gaps and generate hypotheses that are not only innovative but also scientifically valid. This requires balancing creativity and rigor, which is difficult for AI models that rely on existing data and patterns. Additionally, ensuring the novelty of proposals often demands a deep understanding of the broader scientific context, which can be challenging for models limited by the quality and scope of their training data.

    \item Second, implementing closed-loop feedback for end-to-end experimental validation is another major hurdle. Autonomous systems need to design experiments, execute them, analyze results, and iteratively refine their hypotheses in a seamless loop. This requires integration across multiple domains, such as robotics for experiment execution and advanced analytics for result interpretation. Furthermore, real-world experiments often come with unexpected variables and noise, making it challenging for autonomous systems to adapt and learn effectively. Achieving a truly closed-loop system demands robust coordination, adaptability, and the ability to handle uncertainty, which remain technical and conceptual barriers.

\end{itemize}

To further facilitate the advancement of ASR, we propose the InternAgent, an end-to-end auto-research pipeline, which covers four main modules: self-evolving idea generation, human-interactive feedback, idea-to-methodology construction, and multi-round experiment planning and execution. With the help of the Self-Evolving Human-interactive Idea Generation and Idea-to-Methodology Construction, InternAgent can transform a rough proposal into a detailed and easily implementable method, which further increases the success rate of code implementation process and enhances the efficiency of closed-loop experiments. Besides, by leveraging multi-round experiment planning and execution, InternAgent is capable of designing experimental plans and decomposing the experimental process according to the InternAgent-proposed modules, thereby validating the effectiveness of each InternAgent-generated module through experimentation. 

As shown in Fig.~\ref{fig:figure1}, InternAgent has been validated across \textbf{12} scientific research tasks, and we are excited to see that the experimental results demonstrate the significant value of InternAgent in the entire process from hypothesis generation to experimental validation. For instance, in the Reaction Yield Prediction task, the baseline model only achieved a performance of $24.2\%\pm4.2$, while our model improved it to $34.8\%\pm1.1$ in \textbf{just 12 hours}. In contrast, human researchers typically require \textbf{several months} to achieve a similar level of performance improvement. Another example of performance improvement is the Enhancer Activity Prediction task. The baseline model, DeepSTARR, achieved a result of 0.65. By utilizing InternAgent to search relevant domain literature, automatically generate code, and conduct validation, the performance can be improved to 0.79, representing a promising enhancement. In addition, InternAgent also supports complex project-level modifications and debugging, which consist of multiple code files. These results clearly indicate that InternAgent can autonomously generate ideas and design algorithms, effectively reducing the dependence on human effort in scientific research. To facilitate reproducibility, we have open-sourced both the baselines and the codes generated by InternAgent used in all involved scientific research tasks at \texttt{\url{https://github.com/Alpha-Innovator/InternAgent}}.

Furthermore, the contributions of this paper are summarized below:

\begin{itemize}
    \item \textbf{Unified Multi-agent Framework for Diverse Scientific Research Tasks}: 
    We present InternAgent, a unified closed-loop scientific research framework that can automate the entire research cycle, including idea generation, idea-to-methodology transformation, experiment execution, and result feedback. This unified framework can be directly applied to various scientific research scenarios and fields.
    \item \textbf{Interactive Interfaces for Enhanced Cooperative Research}: InternAgent offers interactive interfaces for human-machine collaboration within the idea generation module and across the entire system. By selecting collaboration modes, such as leveraging AI or human experts, it provides evaluations of idea generation effectiveness and facilitates the assessment, reflection, and documentation of experimental results.
    \item \textbf{Comprehensive Experimental Validation and Human Studies}: We conducted extensive human studies centered around InternAgent, including inviting domain experts to evaluate and score the novelty of ideas generated by InternAgent, and comparing the research efficiency between human researchers and InternAgent. These experiments and human studies are crucial for gaining insights into the capabilities of multi-agent systems in conducting scientific research tasks in open-ended environments. We observed many promising phenomena, while also identifying certain technical modules that require improvement.
\end{itemize}


